\section{相关工作}
\label{sec:related-work}

\subsection{标注稀缺条件下的器官分割}

早期田间作物分割通常依赖颜色指数或卷积编码器--解码器。DeepLabV3+ 通过空洞卷积和低层特征融合兼顾上下文与边界 \citep{chen2018deeplabv3plus}，SegFormer 则使用层次化 Transformer 和轻量解码器获得更强的全局建模能力 \citep{xie2021segformer}。GWFSS 数据集上的比较表明，SegFormer 在随机划分和区域划分下均优于 DeepLabV3+，但茎秆仍明显弱于其他器官 \citep{wang2025gwfss}。Mask2Former 将语义分割统一为类别查询与掩码预测问题 \citep{cheng2022mask2former}，ViT-Adapter 则使预训练视觉 Transformer 能够输出适合密集预测的多尺度特征 \citep{chen2023vitadapter}。

有限标注使单纯监督训练容易过拟合。Mean Teacher 使用学生参数的指数滑动平均构建教师 \citep{tarvainen2017mean}，Guided Distillation 进一步在学生预热阶段同时使用有标注与无标注数据 \citep{berrada2024guided}。弱到强一致性方法通过教师的弱扰动预测监督学生的强扰动视图，已成为半监督语义分割的重要范式 \citep{yang2023unimatch}。ST++ 使用跨训练阶段的预测稳定性选择可靠图像 \citep{yang2022stpp}，U2PL 则将不可靠像素作为负类监督以减少无标注信息浪费 \citep{wang2022u2pl}。在小麦器官分割中，已有方法进一步使用器官比例或特征相似度选择无标注图像，并利用不确定性过滤伪标签 \citep{cao2025detailminority,ghosh2025decon}。这些方法提高了无标注数据利用率，但仍将茎秆伪标签主要视为普通区域，没有显式区分中心线结构与宽度边界的可靠性。

\subsection{细长器官的结构完整性}

茎秆具有占比低、宽度小、遮挡频繁和方向连续等特点。低分辨率特征与常规插值容易模糊其边界，PointRend 因而将不确定边界位置视为需要进一步渲染的点 \citep{kirillov2020pointrend}。这类局部细化能够恢复边界细节，但没有直接约束细长目标的连续性。

拓扑保持损失为管状结构分割提供了另一种思路。soft-clDice 通过预测掩码与参考掩码骨架之间的重合关系度量拓扑精度和拓扑召回率 \citep{shit2021cldice}；中心线交叉熵进一步讨论了拓扑一致性与区域精度之间的平衡 \citep{acebes2024clce}。近期工作还从深浅层特征融合和生长--抑制平衡角度增强不同宽度管状结构的识别 \citep{huang2025harmonyseg}。然而，小麦茎秆并不等同于完整可见的血管或道路，其真实标注可能因叶片遮挡而中断。直接对所有预测强制连通会产生不符合可见区域定义的虚假连接。因此，本文只在多视图一致的可见中心线上施加拓扑监督，并将不可靠边界保留为忽略区域。

\subsection{未知采集域中的器官语义稳定性}

田间图像的域差异同时来自成像设备、光照、背景、基因型和生育阶段。GWFSS 的特征可视化呈现出明显的机构聚类，区域划分下的性能下降也说明随机划分难以反映真实部署难度 \citep{wang2025gwfss}。域泛化语义分割要求训练阶段不访问目标域，学习能够迁移到未知环境的表示 \citep{schwonberg2025dgsurvey}。类别记忆方法通过元学习保存跨域共享的类别模式 \citep{kim2022pinmemory}，多分辨率域适应方法则表明全局上下文与高分辨率局部信息对小目标和跨域鲁棒性均有价值 \citep{hoyer2022hrda}。

对小麦器官而言，不同类别的域稳定区域并不相同。叶片和背景具有较大的内部区域，而茎秆最可靠的位置通常位于细长结构的中心。现有类别原型通常从全部预测像素或置信区域求均值，仍可能受到边界混淆和伪标签噪声影响；在半监督早期把人工标注与伪标注写入同一记忆，还会使错误类别中心持续影响后续特征。本文从拓扑核心而非完整预测区域构建域级原型，以课程式双记忆隔离两类证据，并在读取参照时对每个样本排除其自身域，以减小未知地域上的语义漂移。

\subsection{高分辨率细化的精度与效率}

将图像放大后进行滑窗预测能够提升小目标和细结构的可见尺度。SAHI 通过切片推理增强小目标检测 \citep{akyon2022sahi}，高分辨率域适应方法利用低分辨率全局图像和高分辨率局部裁剪融合上下文与细节 \citep{hoyer2022hrda}。GWFSS 上的多尺度滑窗策略证明了放大输入对茎秆分割的有效性，但密集覆盖多个尺度会产生较高推理开销 \citep{cao2025detailminority}。SparseRefine 使用预测熵选择少量高分辨率像素进行稀疏修正 \citep{liu2024sparserefine}，说明选择性计算是提高高分辨率分割效率的可行方向。

单纯的熵能够定位分类不确定区域，却不能区分普通纹理边界与影响茎秆完整性的结构失效。本文进一步联合茎秆概率、异常骨架端点和跨尺度断裂程度评价候选窗口，使高分辨率计算优先服务于细长器官的真实困难区域。
